% \addchap keeps this chapter unnumbered (in the TOC, but without a chapter
% number), so that Chapter 1 below is Assignment 1, Chapter 2 is Assignment 2,
% etc. -- matching the assignment brief's own numbering, e.g. Section 2.4 in
% this report is Question 2.4 in the assignment.
\addchap{Introduction and Assumptions}
\label{ch:introduction}

\workedby{\tbd{naam}}

\tbd{Korte inleiding: waar gaat dit verslag over, welke opgave, welke editie van
het vak, en hoe is het verslag opgebouwd (Part~1 = Assignments 1--3,
Part~2 = Assignments 4--7; sectienummer = vraagnummer uit de opgave).}

\section{Vessel and idealised geometry}
\label{sec:geometry}

\tbd{Beschrijf het schip en de geidealiseerde doorsneden uit de opgave:
secties A (achterschip/machinekamer), B (parallel midscheeps, drie blokken) en
C (voorschip), met blokindeling en hoofdafmetingen. Voeg hier de figuur of
schets toe waarnaar de latere hoofdstukken verwijzen.}

% \begin{figure}[htb]
%   \centering
%   \includegraphics[width=0.85\textwidth]{figures/<bestandsnaam>}
%   \caption{\tbd{onderschrift}}
%   \label{fig:geometry}
% \end{figure}

\section{Notation and symbols}
\label{sec:notation}

\tbd{Lijst met de symbolen die door het hele verslag gebruikt worden
($a$, $b$, $t_p$, $M_t$, $\theta$, $I_p$, $I_w$, $\omega(s)$, $S_\omega(s)$,
$\tau$, $\sigma_{xx}$, \ldots), zodat de notatie in alle hoofdstukken gelijk is.
Spreek dit samen af voordat je begint met rekenen.}

\section{List of assumptions}
\label{sec:assumptions}

All assumptions made in this report are collected here and referred to from the
individual chapters as A1, A2, \ldots

\begin{enumerate}[label=A\arabic*., leftmargin=*]
  \item \tbd{aanname 1 (bijv.\ dunwandige theorie, lineair elastisch materiaal)}
  \item \tbd{aanname 2}
  \item \tbd{\ldots\ vul aan tijdens het uitwerken; elke aanname die je in een
        hoofdstuk maakt, komt hier ook te staan}
\end{enumerate}

\section{Course objectives and general instructions}
\label{sec:instructions}

\tbd{Controleer onderstaande punten met de opgave van dit jaar en pas ze aan
waar nodig.}

\begin{itemize}
  \item Course objectives: modelling of free and constrained warping torsion of
        solid and hollow beam-like members (stiffness, strength, displacements,
        shear and normal stresses); the generalised beam torsion equation;
        identification of torsion load-carrying members and their efficiency at
        local and global level; calculation of the torsional response of marine
        structures; measurement and monitoring of the torsional response.
  \item Answers are presented in simplified form and, where numerical, rounded
        to three significant digits at the end of each question.
  \item Page limit: \tbd{aantal} pages for the main report (front matter and
        appendices excluded).
  \item Deadlines: Part~1 \tbd{datum}, Part~2 \tbd{datum};
        oral examinations \tbd{data}; consultation hours \tbd{tijd en plaats}.
\end{itemize}
